% !LW recipe = latexmk (xelatex)
\documentclass[aspectratio=169,10pt]{beamer}
% !TEX root = main.tex
\usepackage{fix-cm}
\usepackage{fontspec}
\setsansfont{TeX Gyre Heros}
\usepackage[spanish,es-nodecimaldot]{babel}
\usepackage{booktabs,tabularx,array,tikz,graphicx}
% Estilo compartido por Beamer y las figuras independientes.
\usetikzlibrary{arrows.meta,positioning,calc,babel}
% Colores sRGB obtenidos del PNG oficial de IRYCIS (original en Adobe RGB).
\definecolor{irycisblue}{HTML}{00547A}
\definecolor{irycisgreen}{HTML}{6BBB90}
\colorlet{navy}{irycisblue}
\colorlet{teal}{irycisblue}
\definecolor{muted}{HTML}{566570}
\colorlet{pale}{irycisblue!7!white}
\colorlet{amber}{irycisgreen!75!irycisblue}
\colorlet{palewarm}{irycisgreen!16!white}
\tikzset{>=Stealth, every node/.style={text=navy}}

% Figuras originales para la versión visual; no codifican resultados medidos.
\definecolor{pending}{HTML}{AB6A34}
\colorlet{pendingbg}{pending!8!white}
\tikzset{
  vcard/.style={draw=navy!18,fill=pale,rounded corners=3pt,align=center,inner sep=7pt},
  vgood/.style={vcard,fill=palewarm,draw=irycisgreen!65!navy},
  vpending/.style={vcard,fill=pendingbg,draw=pending!45},
  varr/.style={->,line width=1pt,draw=navy!75},
  vlabel/.style={font=\scriptsize,text=muted,align=center},
  vnum/.style={circle,fill=navy,text=white,font=\bfseries,minimum size=7mm,inner sep=0pt},
  pics/keyboard/.style={code={
    \draw[rounded corners=2pt](-.62,-.32) rectangle(.62,.32);
    \foreach \x in {-.42,-.21,0,.21,.42}{\foreach \y in {-.1,.13}{\draw(\x-.05,\y)--(\x+.05,\y);}}
    \draw(-.28,-.23)--(.28,-.23);}},
  pics/chat/.style={code={
    \draw[rounded corners=3pt](-.65,.02)--(-.65,.48)--(.3,.48)--(.3,.02)--(-.23,.02)--(-.43,-.18)--(-.43,.02)--cycle;
    \draw[rounded corners=2pt](.46,.22)--(.67,.22)--(.67,-.32)--(.45,-.32)--(.45,-.51)--(.23,-.32)--(-.15,-.32)--(-.15,-.17);}},
  pics/note/.style={code={
    \draw[rounded corners=2pt](-.4,-.53) rectangle(.4,.53);
    \foreach \y in {.28,.05,-.18}{\draw(-.23,\y)--(.23,\y);}}},
  pics/clock/.style={code={
    \draw(0,0) circle(.51);\draw(0,.35)--(0,0)--(.26,-.1);}},
  pics/shield/.style={code={
    \draw(0,.57)--(.48,.36)--(.4,-.22)--(0,-.59)--(-.4,-.22)--(-.48,.36)--cycle;
    \draw(-.21,.02)--(-.05,-.15)--(.26,.19);}},
  pics/hospital/.style={code={
    \draw(-.57,-.49) rectangle(.57,.49);\draw(-.16,-.49)--(-.16,-.13)--(.16,-.13)--(.16,-.49);
    \draw[line width=2pt](-.18,.24)--(.18,.24);\draw[line width=2pt](0,.06)--(0,.42);}},
  pics/server/.style={code={
    \foreach \y in {-.32,.04,.4}{\draw[rounded corners=2pt](-.52,\y-.13) rectangle(.52,\y+.13);\fill(-.35,\y)circle(.035);\draw(-.08,\y)--(.33,\y);}}},
  pics/check/.style={code={\draw[line width=2pt](-.36,0)--(-.1,-.26)--(.43,.34);}},
  pics/plug/.style={code={
    \draw(-.26,.05) rectangle(.26,.4);\draw(-.14,.4)--(-.14,.62);\draw(.14,.4)--(.14,.62);\draw(0,.05)--(0,-.2)--(.35,-.2);}},
}
\newenvironment{visual}[1][4.7]{\begin{center}\begin{tikzpicture}[x=1cm,y=1cm,every node/.append style={font=\small,align=center}]\path[use as bounding box](0,0) rectangle(14,#1);}{\end{tikzpicture}\end{center}}

\setbeamercolor{normal text}{fg=navy,bg=white}
\setbeamercolor{structure}{fg=teal}
\setbeamercolor{section in toc}{fg=navy}
\setbeamercolor{frametitle}{fg=navy}
\setbeamercolor{block title}{fg=navy,bg=pale}
\setbeamercolor{block body}{fg=navy,bg=pale}
\setbeamerfont{frametitle}{size=\Large,series=\bfseries}
\setbeamerfont{section in toc}{size=\normalsize}
\setbeamertemplate{navigation symbols}{}
\setbeamertemplate{itemize items}{\small\textbullet}
\setbeamertemplate{section in toc}[sections numbered]
\setbeamersize{text margin left=8mm,text margin right=8mm}
% Archivo oficial sin modificar; trim elimina solo sus márgenes transparentes.
\newcommand{\irycislogo}[1]{\includegraphics[width=#1,trim={20bp 36bp 21bp 49bp},clip]{assets/logo-irycis-oficial.png}}
\newcommand{\currentpart}{\ifnum\value{section}>0\relax\insertsectionnumber/5 · \insertsectionhead\else Introducción · recorrido\fi}
\setbeamertemplate{frametitle}{%
  \vspace{2mm}\strut\insertframetitle\strut\par\vspace{1mm}}
\setbeamertemplate{footline}{%
  \begin{tikzpicture}[remember picture,overlay]
    \ifnum\value{framenumber}>1\relax
      \node[anchor=south,inner sep=0pt] at ([yshift=2mm]current page.south) {\irycislogo{20mm}};
    \fi
  \end{tikzpicture}%
  \hspace{8mm}\textcolor{irycisgreen}{\rule{.6mm}{2mm}}\hspace{1.5mm}%
  \textcolor{muted}{\fontsize{7}{9}\selectfont\currentpart}%
  \hfill\textcolor{muted}{\fontsize{7}{9}\selectfont 12/10/2026 · \insertframenumber/\inserttotalframenumber}%
  \hspace{8mm}\vspace{3mm}}
\newcommand{\topic}[1]{{\color{muted}\small #1}\par\vspace{2mm}}
\newcommand{\sources}[1]{\vfill\vspace{.5mm}{\fontsize{6.8}{8}\selectfont\color{muted}Fuentes: #1\par}}
\newcommand{\link}[2]{\href{#1}{\textcolor{teal}{#2}}}
\newcommand{\takeaway}[1]{\par\vspace{2mm}\begin{beamercolorbox}[sep=1.5mm,wd=\textwidth]{block body}\small #1\end{beamercolorbox}}
\newcommand{\stat}[2]{{\fontsize{34}{37}\selectfont\color{teal}\textbf{#1}}\par\vspace{2mm}#2\par}
\newcolumntype{L}[1]{>{\raggedright\arraybackslash}p{#1}}
\renewcommand{\tabularxcolumn}[1]{>{\raggedright\arraybackslash}p{#1}}
\renewcommand{\arraystretch}{1.22}
\hypersetup{pdftitle={Escribas digitales: presentación ejecutiva de 10 minutos},pdfauthor={Documento de trabajo para médicos y gestores},colorlinks=true,linkcolor=navy,urlcolor=teal}

\begin{document}

% 01 · 0:25
\begin{frame}[t]{Escribas digitales}
\topic{Qué son, por qué importan y cómo evaluar su encaje}
\begin{columns}[T,onlytextwidth]
\begin{column}{.47\textwidth}
\vspace{3mm}
{\Large De la consulta a una\par\textbf{nota clínica revisada}.}
\par\vspace{5mm}
Para un hospital con atención en español y \textbf{Dedalus HCIS}, el sistema de historia clínica del hospital.
\par\vspace{4mm}
\textcolor{teal}{\textbf{10 minutos · 5 bloques}}
\end{column}
\begin{column}{.47\textwidth}
\tableofcontents
\par\vspace{6mm}
\irycislogo{44mm}
\end{column}
\end{columns}
\takeaway{Propuesta: comprobar calidad y utilidad local antes de elegir una solución.}
\sources{Informe completo y guion del proyecto, 18/09/2026. Investigación adicional y comprobaciones selectivas: 07/10/2026. Revisión documental.}
\end{frame}

\section{Qué es}

% 02 · 0:55
\begin{frame}[t]{Un asistente para documentar la consulta}
\topic{El médico conversa con el paciente; el escriba prepara la nota.}
\begin{columns}[c,onlytextwidth]
\begin{column}{.57\textwidth}
\centering\includegraphics[width=\linewidth,height=4.65cm,keepaspectratio]{assets/escriba-consulta-v1.png}
\end{column}
\begin{column}{.41\textwidth}
\centering% !TEX root = ../main.tex
% Las etapas internas se explican en la presentación técnica.
\begin{tikzpicture}[x=1cm,y=1cm,
  scribeStep/.style={vcard,text width=4.2cm,minimum height=1.05cm,font=\small},
  arrow/.style={->,line width=1pt,draw=navy!75}]
\path[use as bounding box] (0,0) rectangle (5.7,4.65);
\node[scribeStep] (input) at (2.85,3.9) {\textbf{Consulta + contexto}\\[1mm]\scriptsize Voz · datos clínicos};
\node[scribeStep,fill=palewarm,draw=irycisgreen!65!navy] (draft) at (2.85,2.35) {\textbf{Borrador de nota}\\[1mm]\scriptsize Información organizada};
\node[scribeStep] (doctor) at (2.85,.8) {\textbf{Revisión del médico}\\[1mm]\scriptsize Aprueba y guarda en HCIS};
\draw[arrow] (input.south)--(draft.north);
\draw[arrow] (draft.south)--(doctor.north);
\end{tikzpicture}

\end{column}
\end{columns}
\takeaway{Historia y observaciones, cuando están disponibles. El médico revisa y aprueba.}
\sources{\link{https://www.england.nhs.uk/long-read/guidance-on-the-use-of-ai-enabled-ambient-scribing-products-in-health-and-care-settings/}{NHS England, v3, 2026}; \link{https://docs.nabla.com/core-api/reference/2026-04-24/user/generate-note}{Nabla: transcripción y contexto}. Esquema conceptual; ilustración generada con IA como metáfora del asistente.}
\end{frame}

% 03 · 1:00
\begin{frame}[t]{Ejemplo: de tres fuentes a una nota SOAP}
\topic{Caso ficticio de seguimiento de hipertensión · cada dato conserva su origen.}
\makebox[\textwidth][c]{% !TEX root = ../main.tex
\begin{tikzpicture}[
  x=1cm,y=1cm,
  inputcard/.style={draw=teal!50,fill=pale,rounded corners=2pt,
    text width=45mm,align=left,inner sep=1.5mm,font=\fontsize{8.5}{10}\selectfont},
  soaprow/.style={draw=teal!35,fill=white,rounded corners=2pt,
    text width=65mm,align=left,inner sep=1mm,font=\fontsize{8.5}{10}\selectfont},
  connector/.style={->,draw=teal!70,line width=.8pt}]
\node[inputcard] (conversation) at (0,1.5) {\textbf{Conversación}\\Paciente: «Bien; tensión en casa 135/85».\\Médico: «Hipertensión estable; mantener tratamiento; revisión en 3 meses».};
\node[inputcard,draw=amber!60,fill=palewarm] (chart) at (0,-.4) {\textbf{Historia clínica registrada}\\Hipertensión; losartán 50 mg/día.};
\node[inputcard] (manual) at (0,-1.55) {\textbf{Observación / entrada manual}\\Tensión medida: 138/86 mmHg.};
\node[draw=teal,fill=teal,text=white,rounded corners=2pt,
  text width=13mm,align=center,inner sep=1mm,font=\fontsize{8}{10}\selectfont] (generator) at (3.6,0) {Borrador SOAP};
\draw[connector] (conversation.east) -- (generator.north west);
\draw[connector] (chart.east) -- (generator.west);
\draw[connector] (manual.east) -- (generator.south west);
\node[soaprow] (s) at (8.35,1.7) {\textbf{S · Subjetivo} \hfill {\color{teal}\scriptsize Conversación}\\Se encuentra bien; tensión en casa 135/85.};
\node[soaprow] (o) at (8.35,.6) {\textbf{O · Objetivo} \hfill {\color{amber}\scriptsize Observación + historia}\\Tensión medida: 138/86 mmHg.\\Historia: losartán 50 mg/día.};
\node[soaprow] (a) at (8.35,-.55) {\textbf{A · Valoración} \hfill {\color{teal}\scriptsize Médico · conversación}\\Hipertensión estable, según el médico.};
\node[soaprow] (p) at (8.35,-1.55) {\textbf{P · Plan} \hfill {\color{teal}\scriptsize Médico · conversación}\\Mantener tratamiento; revisión en 3 meses.};
\draw[connector] (generator.east) -- ++(.35,0) |- (s.west);
\draw[connector] (generator.east) -- ++(.35,0) |- (o.west);
\draw[connector] (generator.east) -- ++(.35,0) |- (a.west);
\draw[connector] (generator.east) -- ++(.35,0) |- (p.west);
\end{tikzpicture}
}
\takeaway{Valoración y plan expresados por el médico; contexto de HCIS pendiente de confirmar.}
\sources{Ejemplo propio, ilustrativo. \link{https://www.ncbi.nlm.nih.gov/books/NBK482263/}{SOAP: Subjetivo, Objetivo, Valoración (Assessment) y Plan}. La ubicación de medicación puede variar según la plantilla.}
\end{frame}

\section{Por qué importa}

% 04 · 1:05
\begin{frame}[t]{El beneficio que buscamos: menos carga documental}
\topic{Recuperar atención y tiempo a lo largo de la jornada.}
\begin{visual}[4.65]
\foreach \x/\icon/\title/\sub in {2.2/keyboard/Menos teclado/durante la consulta,7/chat/Más atención/al paciente,11.8/note/Menos pendientes/al terminar}{
  \node[vcard,minimum width=4cm,minimum height=2.6cm] at(\x,3.15) {};
  \pic[draw=navy,line width=1.2pt] at(\x,3.75){\icon};
  \node[font=\bfseries] at(\x,2.85){\title};
  \node[vlabel] at(\x,2.4){\sub};
}
\node[vgood,minimum width=13.5cm,minimum height=1.4cm] at(7,.75){};
\node[vlabel,text=navy,font=\scriptsize] at(2.5,1.19){Proyección ilustrativa};
\node[font=\fontsize{27}{30}\selectfont\bfseries] at(2.5,.62){$\approx$5 h/mes};
\node[align=left,anchor=west,text width=8.5cm] at(4.85,.82){\textbf{menos de documentación por profesional}\\[1mm]\scriptsize Base: 16 min menos por 8 h de consulta programada\\\scriptsize Ejemplo: 8 h de consulta/día · 220 días/año};
\end{visual}
\sources{\link{https://jamanetwork.com/journals/jama/fullarticle/2847319}{Rotenstein et al., JAMA, 2026}: asociación observada en cinco centros de EE. UU. Cálculo propio: 4,9 h/mes medio con efecto constante. No reducción demostrada de jornada. Objetivos superiores por evaluar localmente.}
\end{frame}

% 05 · 0:50
\begin{frame}[t]{Hay que medir hasta que la nota queda cerrada}
\topic{El ahorro depende del trabajo que sigue al borrador.}
\begin{visual}[4.65]
\node[vcard,minimum width=2.7cm,minimum height=1.6cm] (draft) at(1.5,3){\textbf{Borrador}\\generado};
\node[vgood,minimum width=2.7cm,minimum height=1.6cm] (review) at(5,3){\textbf{Revisar}\\y corregir};
\node[vgood,minimum width=2.7cm,minimum height=1.6cm] (save) at(8.5,3){\textbf{Incorporar}\\a HCIS};
\node[vcard,minimum width=2.7cm,minimum height=1.6cm] (closed) at(12,3){\textbf{Nota}\\cerrada};
\draw[varr](draft)--(review);\draw[varr](review)--(save);\draw[varr](save)--(closed);
\draw[navy!55,line width=.8pt](3.55,1.95)--(3.55,1.7)--(13.45,1.7)--(13.45,1.95);
\pic[draw=navy,line width=1.1pt,scale=.6] at(4.5,.85){clock};
\node[anchor=west,align=left] at(5.05,.85){\textbf{Tiempo total + correcciones}\\\scriptsize Calidad clínica y uso correcto del contexto};
\end{visual}
\sources{Propuesta local de evaluación. \link{https://jamanetwork.com/journals/jama/fullarticle/2847319}{JAMA, 2026}; \link{https://www.jmir.org/2025/1/e64993/}{Biro et al., 2025}. La comparación con el flujo habitual debe incluir revisión, corrección y cierre.}
\end{frame}

\section{Mercado y encaje}

% 06 · 0:50
\begin{frame}[t]{El mercado de escribas crece con rapidez}
\topic{Ingresos estimados de escribas ambientales · Menlo Ventures, 2025.}
\begin{visual}[4.65]
\node[anchor=west,font=\scriptsize,text=muted] at(.5,4.38){Millones de dólares};
\foreach \y/\value in {.7/0,1.7/200,2.7/400,3.7/600}{
  \draw[navy!10](1.25,\y)--(7.1,\y);
  \node[anchor=east,font=\scriptsize,text=muted] at(1.1,\y){\value};
}
\fill[navy!22,rounded corners=2pt](2,.7)rectangle(3.7,1.95);
\fill[irycisgreen,rounded corners=2pt](4.9,.7)rectangle(6.6,3.7);
\node[font=\bfseries] at(2.85,2.22){250*};
\node[font=\bfseries\Large] at(5.75,3.99){600};
\node at(2.85,.33){2024*};\node at(5.75,.33){2025};
\node[font=\fontsize{44}{48}\selectfont\bfseries] at(10.6,3.05){$\times$2,4};
\node[font=\large] at(10.6,2.17){en ingresos, en un año};
\node[align=center,text width=5.2cm,font=\small,text=navy] at(10.6,1.15){\textbf{PHTI: entre las adopciones más rápidas de la sanidad reciente}\\[1mm]\scriptsize Valoración sectorial · 2025};
\end{visual}
\sources{\link{https://menlovc.com/perspective/2025-the-state-of-ai-in-healthcare/}{Menlo Ventures, 2025}: estimación de ingresos, EE. UU.; fuente inversora. *2024 derivado: 600/2,4. \link{https://www.phti.org/announcement/ai-scribes-reduce-clinician-burnout/}{PHTI, 25/03/2025}: valoración basada en entrevistas y grupo de expertos, sin ranking histórico cuantitativo.}
\end{frame}

% 07 · 0:25
\begin{frame}[t]{Muchas ofertas; varias formas de llegar al hospital}
\topic{Cuatro familias orientativas, que se solapan.}
\begin{visual}[4.65]
\node[font=\fontsize{40}{43}\selectfont\bfseries] at(1.7,3.55){43};
\node[text width=3cm,font=\small] at(1.7,2.5){soluciones y proyectos en nuestro inventario};
\node[vlabel,text width=3cm] at(1.7,1.65){Septiembre de 2026\\Muestra documental};
\node[vlabel,text width=3.2cm,text=navy] at(1.7,.7){Quirónsalud: proyecto propio\\Referencia de implantación};
\foreach \x/\y/\icon/\title/\example in {6.2/3.4/keyboard/Uso individual/Heidi · Noa · Freed,11.1/3.4/hospital/Oferta\\hospitalaria/Tandem · Invox · Align,6.2/1.25/note/Con la historia\\clínica/Microsoft · Oracle · Abridge,11.1/1.25/server/Para construir\\con APIs/Corti · AWS}{
 \node[vcard,minimum width=4.3cm,minimum height=1.8cm] at(\x,\y){};
 \pic[draw=navy,line width=1pt,scale=.55] at(\x-1.45,\y+.25){\icon};
 \node[font=\bfseries\small,text width=3.1cm] at(\x+.3,\y+.27){\title};
 \node[vlabel,text width=4cm] at(\x,\y-.43){\example};
}
\end{visual}
\sources{Inventario del informe: 09/09/2026 + AWS, 17/09. Familias orientativas. \link{https://www.quironsalud.com/en/communication/current-events/quironsalud-redisena-modelo-asistencial-reafirmandose-human}{Quirónsalud}: proyecto implantado en su red; sin oferta de API comercial acreditada aquí. El recuento no es un censo de proveedores comerciales.}
\end{frame}

% 08 · 0:35
\begin{frame}[t]{El encaje que buscamos tiene cuatro piezas}
\topic{Una buena demo debe convertirse en un flujo viable en nuestro hospital.}
\begin{visual}[4.65]
\node[vgood,minimum width=3cm,minimum height=1.1cm,font=\large\bfseries] (fit) at(7,2.4){Encaje local};
\foreach \x/\y/\name/\title/\sub/\icon in {2.1/3.6/data/Datos y\\seguridad/Condiciones verificadas/shield,11.9/3.6/quality/Calidad en español/Notas clínicas fiables/note,2.1/1.2/integration/HCIS/Leer contexto y guardar/hospital,11.9/1.2/cost/Viabilidad/Coste total y soporte/clock}{
 \node[vcard,minimum width=4cm,minimum height=1.65cm] (\name) at(\x,\y){};
 \pic[draw=navy,line width=1pt,scale=.48] at(\x-1.32,\y+.21){\icon};
 \node[font=\bfseries,text width=2.8cm] at(\x+.36,\y+.2){\title};
 \node[vlabel,text width=3.7cm] at(\x,\y-.41){\sub};
 \draw[varr](\name)--(fit);
}
\end{visual}
\sources{Requisitos propuestos en el informe del proyecto: datos, seguridad, calidad, integración y coste. HCIS exige comprobar paciente/episodio, lectura de contexto y devolución de la nota revisada.}
\end{frame}

% 09 · 0:55
\begin{frame}[t]{Hay candidatos; el encaje sigue por demostrar}
\topic{Selección corta del informe · principal incógnita visible: el flujo con HCIS.}
\begin{visual}[4.7]
\foreach \x/\title/\status/\detail in {2.2/Tandem/HCIS no disponible/según su ficha pública,7/Invox Genesis/API disponible/conector HCIS por demostrar,11.8/Solventum Align/Integración HCE/HCIS específico por confirmar}{
 \node[vcard,minimum width=4.25cm,minimum height=3cm] at(\x,3){};
 \node[font=\bfseries\large,text width=4cm] at(\x,4.05){\title};
 \pic[draw=navy,line width=1.2pt,scale=.75] at(\x-.65,3.12){plug};
 \draw[dashed,draw=pending,line width=1.2pt](\x-.35,2.97)--(\x+.1,2.97);
 \node[draw=navy!40,rounded corners=1pt,font=\scriptsize,inner sep=3pt] at(\x+.7,2.97){HCIS};
 \node[font=\bfseries\small,text width=3.95cm] at(\x,2.36){\status};
 \node[vlabel,text width=4cm] at(\x,1.86){\detail};
}
\node[vlabel] at(7,1.15){También por cerrar para esta selección};
\foreach \x/\title in {2.2/Calidad clínica local,7/Expediente de datos,11.8/Coste total}{
 \node[vpending,minimum width=4.2cm,minimum height=.6cm,font=\small] at(\x,.52){\title};
}
\end{visual}
\sources{\link{https://tandemhealth.ai/es/producto/integraciones/hp-his-hp-hcis}{Tandem HCIS}; \link{https://developer.invoxmedical.com/}{Invox API}; \link{https://www.solventum.com/es-es/home/health-information-technology/solutions/fluency-align/}{Align}. Revisión selectiva 09/10/2026; declaraciones comerciales, sin prueba funcional. Falta de evidencia local no implica imposibilidad técnica.}
\end{frame}

\section{Discusión: cómo evaluar}

% 10 · 0:45
\begin{frame}[t]{Dos vías que podemos poner a prueba}
\topic{La vía local es una opción de exploración, sujeta a recursos.}
\begin{visual}[4.65]
\node[vcard,minimum width=5.9cm,minimum height=2.75cm] (vendor) at(3.25,3.1){};
\node[vgood,minimum width=5.9cm,minimum height=2.75cm] (local) at(10.75,3.1){};
\pic[draw=navy,line width=1.1pt,scale=.7] at(3.25,3.85){hospital};
\pic[draw=navy,line width=1.1pt,scale=.7] at(10.75,3.85){server};
\node[font=\bfseries\large] at(3.25,3.05){Proveedor};
\node[font=\bfseries\large] at(10.75,3.05){Prototipo local};
\node[vlabel,text width=5.3cm] at(3.25,2.26){Producto y soporte existentes\\Condiciones + integración por cerrar};
\node[vlabel,text width=5.3cm] at(10.75,2.26){Control de datos y versiones\\Desarrollo + operación propios};
\node[vgood,minimum width=6.3cm,minimum height=.8cm,font=\bfseries] (compare) at(7,.6){Misma evaluación de calidad y esfuerzo};
\draw[varr](vendor.south)--(3.25,1.17)--(compare.west);
\draw[varr](local.south)--(10.75,1.17)--(compare.east);
\end{visual}
\sources{Propuesta del proyecto. Referencias exploratorias: \link{https://github.com/sammargolis/OpenScribe}{OpenScribe}, \link{https://github.com/bloodworks-io/phlox}{Phlox}; sin instalación ni validación local. Un prototipo local también necesita seguridad, mantenimiento e integración.}
\end{frame}

% 11 · 0:55
\begin{frame}[t]{Los mismos casos permiten una comparación justa}
\topic{El banco de casos sirve para decidir qué solución merece un piloto.}
\begin{visual}[4.65]
\node[vcard,minimum width=3.1cm,minimum height=2.35cm] (cases) at(1.65,2.65){\textbf{Mismos casos}\\[3mm]Audio\\Contexto fechado\\Observaciones};
\node[vcard,minimum width=2.3cm,minimum height=.9cm] (a) at(5.5,3.7){Solución A};
\node[vcard,minimum width=2.3cm,minimum height=.9cm] (b) at(5.5,2.65){Solución B};
\node[vgood,minimum width=2.3cm,minimum height=.9cm] (c) at(5.5,1.6){Prototipo};
\foreach \n in {a,b,c}{\draw[varr](cases.east)--(3.65,2.65)|-(\n.west);}
\node[vcard,minimum width=3.15cm,minimum height=2.45cm] (notes) at(9.2,2.65){\textbf{Misma revisión}\\[3mm]Errores y omisiones\\Uso del contexto\\Correcciones};
\foreach \n in {a,b,c}{\draw[varr](\n.east)--(7.15,0 |- \n.east)|-(notes.west);}
\node[vgood,minimum width=2.3cm,minimum height=1.3cm,font=\bfseries] (decision) at(12.55,2.65){Elegir\\piloto};
\draw[varr](notes)--(decision);
\node[vlabel,text width=12cm] at(7,.48){Pruebas repetibles por lotes + revisión clínica de muestras y casos señalados};
\end{visual}
\sources{Diseño propuesto, basado en \link{https://www.jmir.org/2025/1/e64993/}{Biro et al., 2025} y \link{https://aclanthology.org/2025.emnlp-industry.104/}{Zhou et al., 2025}. Casos de ajuste separados de evaluación; revisión clínica. El ahorro real se comprueba después, en el piloto asistencial.}
\end{frame}

\section{Conclusión y siguientes pasos}

% 12 · 0:45
\begin{frame}[t]{Siguiente paso: comparar antes de pilotar}
\topic{Cada etapa deja un resultado que permite decidir la siguiente.}
\begin{visual}[4.65]
\foreach \x/\n/\title/\sub in {1.7/1/Cerrar\\condiciones/Datos · HCIS · coste,5.2/2/Preparar casos/Español + referencia,8.7/3/Comparar notas/Calidad + correcciones,12.2/4/Piloto\\supervisado/Ahorro + uso real}{
 \node[vnum] at(\x,4.05){\n};
 \node[vcard,minimum width=3.1cm,minimum height=1.7cm] (step\n) at(\x,2.6){};
 \node[font=\bfseries,text width=2.9cm] at(\x,2.95){\title};
 \node[vlabel,text width=2.9cm] at(\x,2.27){\sub};
}
\draw[varr](step1)--(step2);\draw[varr](step2)--(step3);\draw[varr](step3)--(step4);
\node[vgood,minimum width=12.9cm,minimum height=.85cm,font=\bfseries] at(7,.8){Decisión del siguiente hito: qué solución pilotar y con qué condiciones};
\end{visual}
\sources{Propuesta del proyecto: requisitos y criterios de éxito por acordar; candidatos y posible comparador local. El piloto procede si la comparación acredita calidad y viabilidad suficientes.}
\end{frame}

% 13 · 0:35
\begin{frame}[t]{La decisión ahora: iniciar una evaluación local}
\topic{La oportunidad está clara. El encaje debe demostrarse en nuestro flujo.}
\begin{visual}[4.65]
\node[font=\fontsize{24}{28}\selectfont\bfseries,align=left,anchor=west,text width=13.4cm] at(.3,3.65){Elegir qué pilotar\\con evidencia propia.};
\foreach \x/\icon/\title in {2.3/note/Notas fiables,7/clock/Menos trabajo,11.7/hospital/Integración HCIS}{
 \pic[draw=navy,line width=1.25pt,scale=.6] at(\x,1.95){\icon};
 \node[font=\bfseries] at(\x,1.14){\title};
}
\node[vgood,minimum width=13.3cm,minimum height=.65cm,font=\small] at(7,.25){Para arrancar: un servicio · responsables clínico y técnico · recursos acotados};
\end{visual}
\sources{Conclusión y propuesta del proyecto. Próximo entregable: comparación de calidad, esfuerzo e integración para decidir el alcance del piloto.}
\end{frame}

\end{document}
